\documentclass[10pt]{article}
\usepackage[margin=0.18in]{geometry}
\usepackage{multicol,amsmath,amssymb,mathtools,enumitem,booktabs,array}
\usepackage[T1]{fontenc}
\usepackage{lmodern}
\pagestyle{empty}
\setlength{\parindent}{0pt}
\setlength{\parskip}{0.25pt}
\setlength{\columnsep}{0.12in}
\setlength{\multicolsep}{0pt}
\setlist[itemize]{leftmargin=*,nosep,topsep=0pt}
\newcommand{\sect}[1]{\vspace{0.9pt}{\fontsize{7.2}{7.5}\selectfont\textbf{#1}}\vspace{0.25pt}\hrule\vspace{0.8pt}}
\newcommand{\sub}[1]{\textbf{#1}}
\newcommand{\R}{\mathbb{R}}
\newcommand{\C}{\mathbb{C}}
\newcommand{\rank}{\operatorname{rank}}
\newcommand{\range}{\operatorname{range}}
\newcommand{\nullsp}{\operatorname{null}}
\newcommand{\spn}{\operatorname{span}}
\newcommand{\diag}{\operatorname{diag}}
\newcommand{\tr}{\operatorname{tr}}
\newcommand{\argmin}{\operatorname*{arg\,min}}
\newcommand{\E}{\mathbb{E}}
\newcommand{\Var}{\operatorname{var}}
\newcommand{\Cov}{\operatorname{Cov}}
\newcommand{\Prob}{\mathbb{P}}
\newcommand{\one}{\mathbf{1}}
\begin{document}
\fontsize{6.25}{6.63}\selectfont
\begin{multicols*}{3}

\sect{Linear Maps and Models (Lectures 1--4)}
\sub{Matrix map:} \(A\in\R^{m\times n}\) maps inputs \(x\in\R^n\) to
outputs \(y=Ax\in\R^m\). With columns \(a_j\) and rows \(r_i^T\),
\[
y=\sum_{j=1}^n x_ja_j,\qquad y_i=r_i^Tx,\qquad A_{ij}=e_i^TAe_j
\]
Column \(j\) is the response to unit input \(e_j\); row \(i\) is the
sensitivity pattern for output/sensor \(i\).

\sub{Products / composition:} \(ABx=A(Bx)\), so \(B\) acts first.
Column \(j\) of \(AB\) is \(A\) times column \(j\) of \(B\); row \(i\) of
\(AB\) is row \(i\) of \(A\) times \(B\); and
\[
(AB)_{ij}=r_i^Tb_j
\]
where \(b_j\) is column \(j\) of \(B\). \((AB)^T=B^TA^T\). Left
multiplication mixes rows; right multiplication
mixes columns. A permutation matrix has one \(1\) per row/column,
\(P^{-1}=P^T\), and reorders rows or columns.

\sub{Linear map:} \(f:\R^n\to\R^m\) is linear iff
\[
f(\alpha u+\beta v)=\alpha f(u)+\beta f(v)
\]
for all scalars. Necessarily \(f(0)=0\), and every linear map has a unique
matrix \(A=[f(e_1)\cdots f(e_n)]\). To disprove linearity, one failed test
at \(0\), additivity, or homogeneity suffices.

\sub{Affine map:} \(f(x)=Ax+b\) iff it preserves affine combinations
\(\alpha+\beta=1\):
\[
f(\alpha u+\beta v)=\alpha f(u)+\beta f(v)
\]
Subtract the offset: \(g(x)=f(x)-f(0)\) is linear and \(b=f(0)\).

\sub{Linearization:} for differentiable \(f:\R^n\to\R^m\),
\[
f(x_0+\delta x)\approx f(x_0)+Df(x_0)\delta x,\qquad
(Df)_{ij}=\frac{\partial f_i}{\partial x_j}(x_0)
\]
Thus small deviations satisfy \(\delta y\approx Df(x_0)\delta x\).
For scalar \(f\), \(df\approx\nabla f(x)^Tdx\), and
\[
\nabla(x^TAx)=(A+A^T)x
\]

\sub{Building engineering models:} write one output equation per row and
one unknown/input per column. Check dimensions and units. In estimation,
rows are sensors and columns are parameter signatures; in control, columns
are actuator responses and \(\range(A)\) is the set of achievable outputs.

\sub{Signal-flow / graph matrices:} in a directed linear signal graph,
\(A_{ij}\) is the sum of products of edge gains over all paths from input
\(j\) to output \(i\). Cascading the same map gives \(A^2,A^3,\ldots\).
For adjacency matrix \(G\),
\[
(G^k)_{ij}=\#\{\text{walks of length }k\text{ from }j\text{ to }i\}
\]
with the orientation determined by the chosen adjacency convention.

\sect{Linear Algebra (Lecture 5)}
\sub{Norms:}
\[
\|x\|_p=\left(\sum_i|x_i|^p\right)^{1/p},\qquad
\|x\|_\infty=\max_i|x_i|
\]
\(\|x\|\) means \(\|x\|_2\). It is nonnegative with equality iff \(x=0\),
\(\|\alpha x\|=|\alpha|\|x\|\), and
\(\|x+y\|\le\|x\|+\|y\|\).

\sub{Inner product:} \(\langle u,v\rangle=u^Tv\), or
\(u^TWv\) for \(W=W^T\succ0\). It is symmetric, linear in each real
argument, and positive definite. Its induced norm is
\(\|u\|=\sqrt{\langle u,u\rangle}\). Key identities:
\[
|u^Tv|\le\|u\|\|v\|,\qquad
u^Tv=0\Rightarrow\|u+v\|^2=\|u\|^2+\|v\|^2
\]
\[
\cos\theta=\frac{u^Tv}{\|u\|\|v\|}
\]
Positive inner product means acute (\(<90^\circ\)); zero means right;
negative means obtuse (\(>90^\circ\)).

\sub{Hyperplane / halfspace:}
\[
H=\{x:c^Tx=d\},\qquad
H_\pm=\{x:c^Tx\gtreqless d\}
\]
\(c\ne0\) is normal to the hyperplane, whose dimension is \(n-1\).
Points closer to \(a\) than \(b\) form the halfspace
\[
\|x-a\|\le\|x-b\|
\iff
(b-a)^Tx\le\frac{\|b\|^2-\|a\|^2}{2}
\]
with perpendicular-bisector boundary through \((a+b)/2\).

\sub{Subspace / span:} \(S\subseteq\R^n\) is a subspace iff \(0\in S\)
and \(\alpha u+\beta v\in S\) for all \(u,v\in S\). The span is
\[
\spn(v_1,\ldots,v_k)=\left\{\sum_i\alpha_iv_i\right\}
\]
Sums and intersections of subspaces are subspaces. \(S\cup T\) is a
subspace iff \(S\subseteq T\) or \(T\subseteq S\).

\sub{Independence / basis / dimension:}
\[
\sum_i\alpha_iv_i=0\Rightarrow\alpha_i=0
\]
defines independence. A basis of \(S\) is an independent set that spans
\(S\); every \(x\in S\) then has unique coordinates. \(\dim S\) is the
number of vectors in any basis. In a \(d\)-dimensional space, \(d\)
independent vectors automatically form a basis, as do \(d\) spanning
vectors.

\sub{Coordinates / similarity:} for basis matrix
\(B=[b_1\cdots b_n]\),
\[
x=B\tilde x,\qquad \tilde x=B^{-1}x
\]
If \(y=Ax\), its matrix in \(B\)-coordinates is
\[
\tilde A=B^{-1}AB
\]
Similar matrices describe the same map in different coordinates and have
the same rank, determinant, trace, and eigenvalues.

\sub{Invertibility:} for square \(A\), the following are equivalent:
columns independent; columns span \(\R^n\); \(\nullsp(A)=\{0\}\);
\(\range(A)=\R^n\); every \(Ax=b\) has one solution; \(\det(A)\ne0\);
\(A^{-1}\) exists. Equivalently, \(\det(A)=0\) iff \(A\) is singular.
Useful rules:
\[
(AB)^{-1}=B^{-1}A^{-1},\qquad (A^T)^{-1}=(A^{-1})^T
\]

\sect{Range and Nullspace (Lecture 6)}
\sub{Definitions:}
\[
\range(A)=\{Ax:x\in\R^n\}=\spn\{\text{columns of }A\}
\]
\[
\nullsp(A)=\{x:Ax=0\}
\]
\(Ax=b\) is consistent iff \(b\in\range(A)\). If \(z\in\nullsp(A)\),
then \(A(x+z)=Ax\): nullspace directions are input ambiguity or control
freedom.

\sub{One-to-one / left inverse:} the following are equivalent:
\[
\begin{aligned}
\nullsp(A)=\{0\}&\iff A\text{ is one-to-one}\\
&\iff\text{columns independent}\\
&\iff\exists L:\ LA=I
\end{aligned}
\]
For \(A\in\R^{m\times n}\) with independent columns,
\[
L=(A^TA)^{-1}A^T
\]
because \(\nullsp(A^TA)=\nullsp(A)\). Conversely,
\(LA=I,\ Ax=0\Rightarrow x=LAx=0\).

\sub{Onto / right inverse:} the following are equivalent:
\[
\begin{aligned}
\range(A)=\R^m&\iff A\text{ is onto}\\
&\iff\text{rows independent}\\
&\iff\exists R:\ AR=I
\end{aligned}
\]
For independent rows,
\[
R=A^T(AA^T)^{-1}
\]
Conversely \(AR=I\) means \(x=Rb\) reaches every \(b\).

\sub{Exact solution set:} if \(Ax=b\) is consistent,
\[
x=x_p+z,\qquad z\in\nullsp(A)
\]
It is unique iff \(\nullsp(A)=\{0\}\), otherwise infinite. A wide matrix
has nonzero nullspace, so any consistent system has infinitely many
solutions. A tall full-column-rank matrix has at most one solution, but
some \(b\) are unattainable.

\sub{Fundamental orthogonality:}
\[
\range(A)^\perp=\nullsp(A^T),\qquad
\range(A^T)^\perp=\nullsp(A)
\]
\[
\begin{gathered}
\R^m=\range(A)\oplus\nullsp(A^T),\\
\R^n=\range(A^T)\oplus\nullsp(A)
\end{gathered}
\]

\sect{Rank (Lecture 7)}
\sub{Definition / rank-nullity:}
\[
\begin{gathered}
\rank(A)=\dim\range(A)=\rank(A^T),\\
\rank(A)+\dim\nullsp(A)=n
\end{gathered}
\]
Rank is the maximum number of independent columns or rows and the minimum
intermediate dimension in a factorization \(A=BC\). Full rank means
\(\rank(A)=\min(m,n)\).

\sub{Shape consequences:} tall full rank means independent columns,
one-to-one, left invertible. Wide full rank means independent rows, onto,
right invertible. Square full rank means invertible.

\sub{Product bounds:} if \(B\in\R^{m\times p}\),
\(C\in\R^{p\times n}\),
\[
\begin{gathered}
\rank(B)+\rank(C)-p\le\rank(BC),\\
\rank(BC)\le\min\{\rank(B),\rank(C)\}
\end{gathered}
\]
Invertible pre/post multiplication preserves rank.

\sub{Block identities:}
\[
\rank([A\ B])=\dim(\range(A)+\range(B))
\]
\[
\nullsp\!\begin{bmatrix}A\\B\end{bmatrix}
=\nullsp(A)\cap\nullsp(B)
\]
\[
\rank\diag(A,B)=\rank(A)+\rank(B)
\]
\(\rank([A\ B])=\rank(B)\) iff
\(\range(A)\subseteq\range(B)\).

\sub{Range intersection:} if
\(Z=[Z_A^T\ Z_B^T]^T\) spans \(\nullsp([A\ {-B}])\), then
\(\range(AZ_A)=\range(A)\cap\range(B)\). Retain independent columns of
\(AZ_A\), e.g. by QR.

\sub{Dimension identities:}
\[
\dim(S+T)=\dim S+\dim T-\dim(S\cap T)
\]
\[
(S+T)^\perp=S^\perp\cap T^\perp,\qquad
(S\cap T)^\perp=S^\perp+T^\perp
\]
\(S\oplus T\) means \(S+T\) with unique decomposition, equivalently
\(S\cap T=\{0\}\).

\sect{Orthogonality and QR (Lectures 8--9)}
\sub{Orthonormal columns:} \(U=[u_1\cdots u_k]\) has orthonormal columns
iff \(U^TU=I_k\). Then
\[
\|Ux\|=\|x\|,\qquad U^Ty=(u_1^Ty,\ldots,u_k^Ty)
\]
If \(U\) is square, \(UU^T=I\) and \(U^{-1}=U^T\). If \(U\) is tall,
\(UU^T\ne I\) in general.

\sub{\(2\times2\) orthogonal maps:} \(\det U=+1\) means a rotation;
\(\det U=-1\) means a reflection. Orthogonal products and inverses are
orthogonal.

\sub{Projection:} \(P=UU^T\) projects orthogonally onto \(\range(U)\):
\[
\begin{gathered}
P=P^T=P^2,\qquad y=Py+(I-P)y,\\
Py\perp(I-P)y
\end{gathered}
\]
\[
\range(I-P)=\nullsp(P)=\range(P)^\perp
\]
\(\|Py\|\le\|y\|\), with equality iff \(y\in\range(P)\). If orthogonal
projectors \(P,Q\) commute, \(PQ\) projects onto
\(\range(P)\cap\range(Q)\).

\sub{Gram--Schmidt:} for independent columns \(a_j\),
\[
\tilde q_j=a_j-\sum_{i<j}(q_i^Ta_j)q_i,\qquad
q_j=\tilde q_j/\|\tilde q_j\|
\]
This preserves every leading span:
\[
\spn(q_1,\ldots,q_j)=\spn(a_1,\ldots,a_j)
\]

\sub{Thin QR:}
\[
A=QR,\qquad Q^TQ=I,\qquad R\text{ upper triangular}
\]
\[
r_{ij}=q_i^Ta_j\ (i<j),\qquad r_{jj}=\|\tilde q_j\|
\]
For independent columns \(R\) is invertible and
\(\range(Q)=\range(A)\). If a residual is zero, skip that dependent
column; the number of generated \(q_i\)'s is \(\rank(A)\).

\sub{Full / pivoted QR:}
\[
A=[Q_1\ Q_2]\begin{bmatrix}R\\0\end{bmatrix},\qquad
\range(Q_2)=\range(A)^\perp
\]
Pivoted QR uses \(AP=QR\) to move independent/informative columns first.
Gram--Schmidt depends on column order. Replacing \(a_j\) by cumulative
sums preserves leading spans and \(Q\) up to signs, but changes \(R\).

\sub{QR solves:} for full-column-rank \(A=QR\),
\[
Ax=b\text{ consistent}\Rightarrow Rx=Q^Tb,\qquad
P_{\range(A)}=QQ^T
\]
The same triangular equation gives least squares when \(b\) is
inconsistent.

\sect{Homework 1--3 Patterns}
\sub{Matrix construction:} translate each output formula into one row.
Downsample \(y_i=x_{2i}\) uses row \(e_{2i}^T\); pair averaging uses
\(\tfrac12(e_{2i-1}+e_{2i})^T\). For path systems, multiply gains along a
path and add gains across parallel paths.

\sub{PSD/null trick:} if \(M=N^TN\),
\[
x^TMx=\|Nx\|^2,\qquad \nullsp(M)=\nullsp(N)
\]
For a sum \(\sum_iN_i^TN_i\), its nullspace is
\(\bigcap_i\nullsp(N_i)\). This proves invertibility by showing the only
vector annihilated by every term is zero.

\sub{Nonnegative-independence trick:} if every \(u_i\ne0\) and
\(u_i^Tu_j\ge0\), then
\[
\sum_i\alpha_iu_i=0,\quad\alpha_i\ge0
\quad\Rightarrow\quad \alpha_i=0\ \forall i
\]
because the squared norm is a sum of nonnegative terms and includes
\(\alpha_i^2\|u_i\|^2\).

\sub{Rational interpolation:}
\[
f(t)=\frac{a_0+\cdots+a_mt^m}{1+b_1t+\cdots+b_mt^m}
\]
Each sample \((t_i,y_i)\) gives the linear row
\[
\begin{bmatrix}
1&t_i&\cdots&t_i^m&-y_it_i&\cdots&-y_it_i^m
\end{bmatrix}\theta=y_i
\]
Increase \(m\) until consistent; then verify residual, actual interpolation
error, and nonzero denominators at all samples.

\sub{Wireless feasibility:} with desired common SINR \(S\), direct gains
\(D=\diag(G_{ii})\), interference \(K=G-D\), and noise \(\sigma\one\),
\[
(D-SK)p=S\sigma\one
\]
For each \(S\), solve once and check \(0\le p_i\le p_i^{\max}\); this is a
feasibility test, not an optimization problem.

\sub{Dynamics / coin collection:} piecewise-constant forces give linear
positions \(x(t)=\sum_j a_{tj}f_j\). Collecting selected targets is
\(Af=x_{\rm target}\); consistency decides feasibility. To identify one
target that must be missed, delete each target row and test the remaining
system.

\sub{Finite fault search:} when each candidate fault determines a model
matrix, solve the linear system for every candidate and compare its
predicted scalar measurement with the observed one. Finite model choices
can be identified even from one scalar when their predictions differ.

\sect{Least Squares (Lecture 10)}
\sub{Problem / residual:} for \(A\in\R^{m\times n}\), usually \(m>n\),
\[
x_{\rm ls}=\argmin_x\|Ax-y\|^2,\qquad r_{\rm ls}=Ax_{\rm ls}-y
\]
Geometrically, \(Ax_{\rm ls}\) is the closest point in \(\range(A)\) to
\(y\). Exact fit occurs iff \(y\in\range(A)\).

\sub{Normal equations:} differentiating gives
\[
A^T(Ax_{\rm ls}-y)=0,\qquad A^TAx_{\rm ls}=A^Ty
\]
Thus \(r_{\rm ls}\perp\range(A)\). If \(A\) has independent columns,
\[
x_{\rm ls}=(A^TA)^{-1}A^Ty=A^\dagger y
\]
The minimizer is unique iff \(\nullsp(A)=\{0\}\). If \(A\) is rank
deficient, the fitted vector \(Ax_{\rm ls}\) is still unique, but the
coefficient vector need not be.

\sub{Projection:} for full column rank \(A\),
\[
P_A=A(A^TA)^{-1}A^T,\qquad \hat y=P_Ay,\qquad r=(P_A-I)y
\]
\(P_A=P_A^T=P_A^2\), \(\range(P_A)=\range(A)\), and
\(\nullsp(P_A)=\nullsp(A^T)\). Also
\[
y=P_Ay+(I-P_A)y,\qquad P_Ay\perp(I-P_A)y
\]

\sub{Completion of squares:} for any \(x\),
\[
\|Ax-y\|^2=\|r_{\rm ls}\|^2+\|A(x-x_{\rm ls})\|^2
\]
because \(r_{\rm ls}\perp A(x-x_{\rm ls})\). This proves optimality and,
when \(A\) has zero nullspace, uniqueness.

\sub{QR solution:} if \(A=QR\) is thin QR with \(Q^TQ=I\) and invertible
\(R\),
\[
x_{\rm ls}=R^{-1}Q^Ty,\qquad \hat y=QQ^Ty,\qquad r=(QQ^T-I)y
\]
Solve \(Rx=Q^Ty\); do not explicitly form \(R^{-1}\) or \(A^TA\).

\sub{Least-squares estimation / BLUE:} in \(y=Ax+v\), a linear estimate
\(\hat x=By\) is unbiased in the noise-free case iff \(BA=I\). Then
\(x-\hat x=-Bv\). For full-column-rank \(A\), \(A^\dagger\) is the
smallest left inverse in Frobenius norm:
\[
BA=I\quad\Longrightarrow\quad \|B\|_F^2\ge\|A^\dagger\|_F^2
\]
With equal, uncorrelated sensor noise this is the best linear unbiased
estimator.

\sub{Weighted least squares:} for positive definite weight \(W\),
\[
\min_x(Ax-y)^TW(Ax-y)
\quad\Rightarrow\quad
x=(A^TWA)^{-1}A^TWy
\]
when \(A^TWA\) is invertible. Equivalently apply ordinary LS to
\(W^{1/2}A,W^{1/2}y\).

\sub{Unknown interference:} for \(y=Ax+Bv+w\), set
\(P^\perp=I-BB^\dagger\). Joint LS over \(x,v\) eliminates \(v\):
\[
\hat x=(P^\perp A)^\dagger P^\perp y
\]
It is identifiable noise-free iff \(P^\perp A\) has independent columns;
for full-column-rank \(A\), iff \(\range(A)\cap\range(B)=\{0\}\).

\sect{Multiobjective LS (Lecture 11)}
\sub{Two objectives:} with
\[
J_1(x)=\|Ax-y\|^2,\qquad J_2(x)=\|Fx-g\|^2
\]
the weighted objective is ordinary least squares:
\[
J_1+\mu J_2=
\left\|
\begin{bmatrix}A\\\sqrt\mu F\end{bmatrix}x-
\begin{bmatrix}y\\\sqrt\mu g\end{bmatrix}
\right\|^2,\qquad \mu\ge0
\]
If \(A^TA+\mu F^TF\) is invertible,
\[
x_\mu=(A^TA+\mu F^TF)^{-1}(A^Ty+\mu F^Tg)
\]
The stacked matrix has independent columns iff
\(\nullsp(A)\cap\nullsp(F)=\{0\}\) for \(\mu>0\).

\sub{Tradeoff curve:} \(x\) is Pareto optimal if no other point improves
one objective without worsening the other. On a plot with horizontal
\(J_2\) and vertical \(J_1\), \(J_1+\mu J_2=\alpha\) has slope \(-\mu\).
\(\mu=0\) prioritizes \(J_1\); \(\mu\to\infty\) prioritizes \(J_2\), then
\(J_1\) among minimizers of \(J_2\).

\sub{Tikhonov / ridge:} with \(F=I,g=0\),
\[
x_\mu=(A^TA+\mu I)^{-1}A^Ty,\qquad \mu>0
\]
\(A^TA+\mu I\succ0\) for every \(A\), since
\[
z^T(A^TA+\mu I)z=\|Az\|^2+\mu\|z\|^2>0\quad(z\ne0)
\]
Large \(\mu\) favors small \(x\); small \(\mu\) favors data fit.

\sub{Push-through identity:}
\[
(A^TA+\mu I)^{-1}A^T=A^T(AA^T+\mu I)^{-1}
\]
It follows by first observing
\((A^TA+\mu I)A^T=A^T(AA^T+\mu I)\). This lets you invert the smaller
matrix.

\sect{Least-Norm Solutions (Lecture 12)}
\sub{Underdetermined equations:} let \(A\in\R^{m\times n}\), \(m<n\),
have full row rank. Every \(y\) is attainable and all solutions are
\[
x=x_p+z,\qquad z\in\nullsp(A)
\]
The unique least-norm solution is
\[
x_{\rm ln}=A^T(AA^T)^{-1}y=A^\dagger y
\]

\sub{Orthogonality proof:} \(x_{\rm ln}\in\range(A^T)=\nullsp(A)^\perp\).
For any solution \(x=x_{\rm ln}+z\), \(z\in\nullsp(A)\), so
\[
\|x\|^2=\|x_{\rm ln}\|^2+\|z\|^2\ge\|x_{\rm ln}\|^2
\]
Also \(A^\dagger\) is a right inverse and
\[
\begin{gathered}
A^\dagger A\text{ projects onto }\range(A^T),\\
I-A^\dagger A\text{ projects onto }\nullsp(A)
\end{gathered}
\]

\sub{Lagrange derivation:} for
\(\min\frac12\|x\|^2\) subject to \(Ax=y\), use
\[
\mathcal L(x,\lambda)=\tfrac12x^Tx+\lambda^T(Ax-y)
\]
Stationarity gives \(x=-A^T\lambda\); feasibility gives
\(\lambda=-(AA^T)^{-1}y\).

\sub{Equality-constrained least squares:}
\[
\min_x\tfrac12\|Ax-b\|^2\quad\text{subject to }Cx=d
\]
has KKT equations
\[
\begin{bmatrix}A^TA&C^T\\C&0\end{bmatrix}
\begin{bmatrix}x\\\lambda\end{bmatrix}
=
\begin{bmatrix}A^Tb\\d\end{bmatrix}
\]
The KKT matrix is invertible iff \(C\) has independent rows and
\(\nullsp(A)\cap\nullsp(C)=\{0\}\). Proof: multiply the first homogeneous
KKT equation by \(x^T\) and use \(Cx=0\) to obtain \(\|Ax\|^2=0\).

\sub{Explicit constrained correction:} if \(H=A^TA\succ0\), let
\(x_{\rm ls}=H^{-1}A^Tb\). Then
\[
x^\star=x_{\rm ls}+H^{-1}C^T(CH^{-1}C^T)^{-1}(d-Cx_{\rm ls})
\]
It is the smallest \(H\)-metric correction that enforces \(Cx=d\).

\sect{Homework 4--5 Patterns}
\sub{Matrix system identification:} with samples as columns,
\(X=[x^{(1)}\cdots x^{(N)}]\), \(Y=[y^{(1)}\cdots y^{(N)}]\),
\[
\min_A\|AX-Y\|_F^2,\qquad
\hat A=YX^T(XX^T)^{-1}
\]
when \(X\) has full row rank. Each row of \(A\) is an independent LS fit.

\sub{Multiplicative model:} for
\(\hat L=\alpha\prod_j z_j^{\beta_j}\), take logs:
\[
\log\hat L=\log\alpha+\sum_j\beta_j\log z_j
\]
Fit \((\log\alpha,\beta_j)\) by LS, then exponentiate. This minimizes
squared log error, approximately relative error for a good fit.

\sub{Tomography / inverse problems:} line integral \(j\) is
\(y_j=\sum_i\ell_{ij}x_i+v_j\), so \(A_{ji}=\ell_{ij}\) and
\(\hat x=A^\dagger y\). Always check residual scale and reshape using the
given pixel ordering.

\sub{Curve smoothing:} second-difference matrix
\[
(Dg)_i=\frac{g_{i+2}-2g_{i+1}+g_i}{h^2}
\]
penalizes curvature. For
\[
\frac1n\|f-g\|^2+\frac{\mu}{n-2}\|Dg\|^2
\]
the solution satisfies
\[
\left(\frac1nI+\frac{\mu}{n-2}D^TD\right)g=\frac1n f
\]
\(\mu=0\Rightarrow g=f\); \(\mu\to\infty\Rightarrow Dg=0\), so \(g\) is
the best affine fit to \(f\).

\sub{Interpolation basis:} if \(B\in\R^{n\times m}\) has independent
columns, select \(m\) independent rows to form invertible \(Z\), then
\[
V=BZ^{-1}
\]
has the same range and identity rows at the selected sample indices.
Consequently every \(y\in\range(B)\) is reconstructed directly from those
\(m\) sampled entries.

\sub{Minimum-energy flow/control:} encode conservation or endpoint
requirements as \(Cf=d\). The minimum-energy input is
\[
f^\star=C^T(CC^T)^{-1}d
\]
for independent constraints. Incidence matrices turn node conservation
into \(Af=-s\).

\sub{Conservation-law correction:} to minimally adjust \(y\) so
\(A^Tx=b\), solve \(A^Tc=b-A^Ty\):
\[
y_{\rm adj}=y+A(A^TA)^{-1}(b-A^Ty)
\]
This is orthogonal projection onto the affine set \(\{x:A^Tx=b\}\).

\sect{Complex Numbers / Complex Linear Algebra}
\sub{Arithmetic:} \(i^2=-1\), \(z=a+ib\), \(\bar z=a-ib\),
\[
|z|=\sqrt{a^2+b^2}=\sqrt{\bar z z},\qquad
z^{-1}=\bar z/|z|^2\quad(z\ne0)
\]
Real-coefficient polynomials, and therefore real matrices, have nonreal
roots/eigenvalues in conjugate pairs.

\sub{Polar form / Euler:}
\[
z=re^{i\theta}=r(\cos\theta+i\sin\theta),\quad
r=|z|,\quad\theta=\arg z
\]
Multiplication multiplies moduli and adds angles:
\[
z_1z_2=r_1r_2e^{i(\theta_1+\theta_2)},\qquad
z^k=r^ke^{ik\theta}
\]
\(|z|<1\) gives decay, \(|z|>1\) growth, and \(|z|=1\) constant magnitude
with rotation/oscillation.

\sub{Complex vectors:} \(u^*=\bar u^T\),
\[
\langle u,v\rangle=u^*v,\qquad \|u\|^2=u^*u
\]
Use Hermitian transpose \(^*\), not plain transpose, in orthogonality,
least squares, and unitary matrices. Complex LS is
\[
x_{\rm ls}=(A^*A)^{-1}A^*y
\]
for full column rank.

\sub{Real representation:} for \(A=A_R+iA_I\), \(x=x_R+ix_I\),
\[
\widetilde A=
\begin{bmatrix}A_R&-A_I\\A_I&A_R\end{bmatrix},\qquad
\widetilde x=\begin{bmatrix}x_R\\x_I\end{bmatrix}
\]
Then \(\widetilde{Ax}=\widetilde A\widetilde x\),
\(\|\widetilde x\|=\|x\|\), and
\(\widetilde{A^*}=\widetilde A^T\).

\sect{Eigenvalues / Diagonalization (Lecture 13)}
\sub{Definitions:}
\[
\begin{gathered}
Av=\lambda v,\quad v\ne0
\quad\Longleftrightarrow\quad (A-\lambda I)v=0,\\
\det(\lambda I-A)=0
\end{gathered}
\]
\(v\) is a right eigenvector. A left eigenvector satisfies
\(w^*A=\lambda w^*\). Eigenvectors may be scaled by any nonzero scalar.

\sub{Action / dynamics:} along \(v\), \(A\) only scales by \(\lambda\).
For \(x_{k+1}=Ax_k\), the mode \(v\) evolves as \(\lambda^kv\). If
\(\lambda=re^{i\theta}\), \(r\) controls growth/decay and \(\theta\)
controls rotation/oscillation.

\sub{Diagonalization:} \(A\in\C^{n\times n}\) is diagonalizable iff it has
\(n\) independent eigenvectors. With
\[
V=[v_1\ \cdots\ v_n],\qquad \Lambda=\diag(\lambda_1,\ldots,\lambda_n)
\]
\[
AV=V\Lambda,\qquad A=V\Lambda V^{-1}
\]
Distinct eigenvalues imply independent eigenvectors and diagonalizability;
the converse is false. A matrix lacking a full eigenbasis is defective.
Rows of \(V^{-1}\) are dual left eigenvectors:
\(w_i^*v_j=\delta_{ij}\).

\sub{Useful consequences:}
\[
A^k=V\Lambda^kV^{-1},\quad
A^{-1}=V\Lambda^{-1}V^{-1}\ (\lambda_i\ne0)
\]
\[
f(A)=V\,\diag(f(\lambda_i))\,V^{-1}
\]
for a polynomial/power-series \(f\). Spectral mapping:
\[
Av=\lambda v\quad\Rightarrow\quad f(A)v=f(\lambda)v
\]
Also \(\tr A=\sum_i\lambda_i\), \(\det A=\prod_i\lambda_i\), counting
algebraic multiplicity.

\(AB\) and \(BA\) have the same nonzero eigenvalues:
\(ABv=\lambda v,\lambda\ne0\Rightarrow Bv\ne0\) and
\(BA(Bv)=\lambda Bv\).

\sub{Matrix roots / logarithms:} for diagonalizable
\(A=V\Lambda V^{-1}\), choose scalar values \(\mu_i^2=\lambda_i\) or
\(e^{\nu_i}=\lambda_i\):
\[
A^{1/2}=V\diag(\mu_i)V^{-1},\qquad
\log A=V\diag(\nu_i)V^{-1}
\]
The logarithm requires \(A\) invertible. Branches make these nonunique.
A real diagonalizable matrix has a real
root/log only when choices can be paired conjugately; each distinct
negative real eigenvalue must have even multiplicity.

\sect{Symmetric Matrices (Lecture 14)}
\sub{Spectral theorem:} if \(A=A^T\in\R^{n\times n}\), all eigenvalues are
real and an orthonormal eigenbasis exists:
\[
A=Q\Lambda Q^T=\sum_{i=1}^n\lambda_iq_iq_i^T,\qquad Q^TQ=I
\]
Eigenvectors for distinct eigenvalues are orthogonal; within a repeated
eigenspace they can be chosen orthonormal.

\sub{Geometric action:} \(Q^Tx\) resolves \(x\) into eigen-coordinates,
\(\Lambda\) scales each coordinate, and \(Q\) reconstructs. Thus a
symmetric map has orthogonal input/output principal directions.

\sub{Quadratic forms:} \(f(x)=x^TAx\) depends only on
\[
A_s=\tfrac12(A+A^T),\qquad x^TAx=x^TA_sx
\]
For symmetric \(A\), \(\|x\|=1\), and
\(\lambda_1\ge\cdots\ge\lambda_n\),
\[
\lambda_n\le x^TAx\le\lambda_1
\]
More generally,
\[
\lambda_{\min}\|x\|^2\le x^TAx\le\lambda_{\max}\|x\|^2
\]
The extrema are attained at corresponding eigenvectors.

\sub{PSD / PD:}
\[
\begin{array}{c|c|c}
 &x^TAx&\text{eigenvalues}\\ \hline
A\succeq0&\ge0\ \forall x&\lambda_i\ge0\\
A\succ0&>0\ \forall x\ne0&\lambda_i>0
\end{array}
\]
\(A\succ0\Rightarrow A\) invertible and \(A^{-1}\succ0\).
\(A\succeq0\) iff \(A=F^TF\) for some \(F\). The minimum number of rows
of \(F\) is \(\rank(A)\).

\sub{Indefinite forms / inertia:} split the spectral decomposition into
positive and negative eigenvalues:
\[
A=Q_+\Lambda_+Q_+^T-Q_-\Lambda_-Q_-^T
\]
Then
\[
x^TAx=\|\Lambda_+^{1/2}Q_+^Tx\|^2-
\|\Lambda_-^{1/2}Q_-^Tx\|^2
\]
The minimum row counts are the numbers \(n_+(A)\), \(n_-(A)\) of positive
and negative eigenvalues.

\sub{Matrix inequalities:} \(A\succeq B\) means \(A-B\succeq0\), or
\(x^TAx\ge x^TBx\) for all \(x\). Useful rules:
\[
\begin{gathered}
A\succeq B\Rightarrow -A\preceq-B,\\
A\succeq B\Rightarrow C^TAC\succeq C^TBC,\\
A\succeq B\succ0\Rightarrow A^{-1}\preceq B^{-1}
\end{gathered}
\]
The order is partial. In general \(A\succeq B\) does not imply
\(A^2\succeq B^2\).

\sub{Symmetric square root:} if \(A\succeq0\),
\[
A^{1/2}=Q\diag(\sqrt{\lambda_i})Q^T,\qquad
A^{1/2}A^{1/2}=A
\]
For \(A\succ0\), \(A^{-1/2}\) exists.

\sect{Ellipsoids (Lecture 15)}
\sub{Quadratic representation:} \(c\in\R^n\) is the center; for
\(S=S^T\succ0\),
\[
\mathcal E=\{x:(x-c)^TS(x-c)\le1\}
\]
If \(Sq_i=\lambda_iq_i\), principal directions are \(q_i\) and semiaxis
lengths are \(1/\sqrt{\lambda_i}\). Large \(\lambda_i\) means a short,
thin direction. Boundary parametrization:
\[
x=c+Q\Lambda^{-1/2}u,\qquad \|u\|=1
\]

\sub{Scaled / covariance ellipsoid:}
\[
(x-c)^T\Sigma^{-1}(x-c)\le\alpha
\]
has semiaxis lengths \(\sqrt{\alpha\lambda_i(\Sigma)}\). Multiplying
\(\alpha\) by four doubles all lengths.

\sub{Image-of-ball representation:}
\[
\mathcal E=\{c+Bz:\|z\|\le1\},\qquad B\text{ invertible}
\]
is equivalent to \(S=(BB^T)^{-1}\). Given \(S\), all valid maps are
\[
B=S^{-1/2}R,\qquad R\in\R^{n\times n},\quad R^TR=I
\]
\(R\) is any rotation or reflection of the unit ball, explaining why
the quadratic matrix \(S\) is unique while \(B\) is not.

\sub{Sensors:} if \(y=B^Tx\), then
\[
\|y\|^2\le1\quad\Longleftrightarrow\quad x^TBB^Tx\le1
\]
This is a bounded ellipsoid iff \(BB^T\succ0\), i.e. the sensor vectors
span the state space.

\sub{Fitting an ellipsoid:} \(x_k^TSx_k\approx1\) is linear in the unique
symmetric entries of \(S\). A design row contains \(x_i^2\) on diagonal
terms and \(2x_ix_j\) on off-diagonal terms. Solve
\[
\Phi\theta\approx\one
\]
then reconstruct \(S\), verify \(S\succ0\), and assess
\((x_k^TSx_k-1)\) using RMS and maximum absolute error.

\sect{Matrix Norm and Gain (Lecture 16)}
\sub{Maximum gain / induced \(2\)-norm:}
\[
\begin{gathered}
\|A\|=\max_{x\ne0}\frac{\|Ax\|_2}{\|x\|_2}
=\max_{\|x\|_2=1}\|Ax\|_2,\\
\|A\|=\sqrt{\lambda_{\max}(A^TA)}=\sigma_{\max}(A)
\end{gathered}
\]
The maximizing input is the top eigenvector of \(A^TA\); its output
direction is the corresponding left singular vector.

\sub{Minimum gain:}
\[
\sigma_{\min}(A)=\min_{\|x\|=1}\|Ax\|
=\sqrt{\lambda_{\min}(A^TA)}
\]
It is zero iff \(\nullsp(A)\ne\{0\}\). For a square invertible matrix it
is the smallest singular value and
\(\|A^{-1}\|=1/\sigma_{\min}(A)\).

\sub{Gain bounds / norm properties:} for \(x\ne0\),
\[
\begin{gathered}
\sigma_{\min}(A)\le\frac{\|Ax\|}{\|x\|}
\le\sigma_{\max}(A)=\|A\|,\\
\|AB\|\le\|A\|\,\|B\|,\quad \|A^T\|=\|A\|,\\
\|\alpha A\|=|\alpha|\|A\|,\quad
\|A+B\|\le\|A\|+\|B\|
\end{gathered}
\]
For square \(X,Y\),
\[
\begin{gathered}
\sigma_{\max}(XY)\le\sigma_{\max}(X)\sigma_{\max}(Y),\\
\sigma_{\min}(XY)\ge\sigma_{\min}(X)\sigma_{\min}(Y),\\
\sigma_{\min}(X+Y)\ge\sigma_{\min}(X)-\sigma_{\max}(Y)
\end{gathered}
\]

\sub{Frobenius norm:}
\[
\|A\|_F=\sqrt{\sum_{ij}A_{ij}^2}
=\sqrt{\tr(A^TA)}
=\sqrt{\sum_i\sigma_i^2}
\]
It is the Euclidean norm of all entries, not the maximum gain.
\[
\|A\|\le\|A\|_F\le\sqrt{\rank(A)}\,\|A\|
\]

\sub{Norm invertibility test:} \(\|E\|<1\) implies
\[
(I-E)^{-1}=\sum_{k\ge0}E^k,\qquad
\|(I-E)^{-1}\|\le(1-\|E\|)^{-1}
\]
Thus \(A+\Delta A\) is invertible if \(\|A^{-1}\Delta A\|<1\).

\sect{SVD (Lecture 17)}
\sub{Thin SVD:} if \(A\in\R^{m\times n}\) has rank \(r\),
\[
A=U_r\Sigma_rV_r^T=\sum_{i=1}^r\sigma_i u_iv_i^T
\]
\[
\begin{gathered}
U_r\in\R^{m\times r},\quad V_r\in\R^{n\times r},\\
\Sigma_r=\diag(\sigma_1,\ldots,\sigma_r),\quad
\sigma_1\ge\cdots\ge\sigma_r>0
\end{gathered}
\]
Columns of \(U_r,V_r\) are orthonormal, and
\[
Av_i=\sigma_i u_i,\qquad A^Tu_i=\sigma_i v_i
\]

\sub{SVD and four spaces:}
\[
\begin{array}{c|c}
\range(A)&\spn\{u_1,\ldots,u_r\}\\
\range(A^T)&\spn\{v_1,\ldots,v_r\}\\
\nullsp(A^T)&\spn\{\text{remaining left }u_i\}\\
\nullsp(A)&\spn\{\text{remaining right }v_i\}
\end{array}
\]
\(A^TA=V_r\Sigma_r^2V_r^T\) and
\(AA^T=U_r\Sigma_r^2U_r^T\): nonzero eigenvalues are \(\sigma_i^2\).

\sub{Full SVD:}
\[
A=U\Sigma V^T,\qquad U\in\R^{m\times m},\quad
V\in\R^{n\times n}
\]
with orthogonal \(U,V\) and rectangular diagonal \(\Sigma\), whose first
\(r\) entries are positive and the rest zero. It means: rotate by \(V^T\),
scale/truncate/zero-pad by \(\Sigma\), then rotate by \(U\).

\sub{Symmetric case:} if \(A=A^T\), its singular values are the absolute
eigenvalues, sorted by magnitude:
\[
\{\sigma_i(A)\}=\{|\lambda_i(A)|\}
\]
For \(\lambda_i\ne0\), one choice is
\(v_i=q_i,\ u_i=\operatorname{sign}(\lambda_i)q_i\).

\sub{Moore--Penrose pseudoinverse:}
\[
A^\dagger=V_r\Sigma_r^{-1}U_r^T
\]
For every \(A,y\), \(A^\dagger y\) is the minimum-norm vector among all
least-squares minimizers. All LS minimizers are
\[
x=A^\dagger y+z,\qquad z\in\nullsp(A)
\]
The four identities are
\[
\begin{gathered}
AA^\dagger A=A,\qquad A^\dagger AA^\dagger=A^\dagger,\\
(AA^\dagger)^T=AA^\dagger,\qquad(A^\dagger A)^T=A^\dagger A
\end{gathered}
\]
Thus
\[
AA^\dagger=P_{\range(A)},\qquad
A^\dagger A=P_{\range(A^T)}
\]
and \(I-AA^\dagger=P_{\nullsp(A^T)}\),
\(I-A^\dagger A=P_{\nullsp(A)}\).

\sub{Regularization limit / filters:}
\[
A^\dagger=\lim_{\mu\downarrow0}(A^TA+\mu I)^{-1}A^T
\]
in the action on data. In singular coordinates,
\[
x_\mu=\sum_{i=1}^r
\frac{\sigma_i}{\sigma_i^2+\mu}(u_i^Ty)v_i
\]
Ordinary LS uses gain \(1/\sigma_i\); Tikhonov multiplies it by
\(\sigma_i^2/(\sigma_i^2+\mu)\), strongly damping small-\(\sigma_i\)
directions.

\sub{Sensitivity / condition number:} for invertible \(A\),
\[
\kappa(A)=\|A\|\|A^{-1}\|
=\frac{\sigma_{\max}}{\sigma_{\min}}\ge1
\]
If \(Ax=y\) and only \(y\) is perturbed,
\[
\frac{\|\delta x\|}{\|x\|}
\le\kappa(A)\frac{\|\delta y\|}{\|y\|}
\]
A small \(\sigma_{\min}\) means large noise amplification and proximity
to singularity.

\sub{Best low-rank approximation:} truncate after \(k<r\):
\[
A_k=\sum_{i=1}^k\sigma_i u_iv_i^T
\]
Eckart--Young:
\[
\begin{gathered}
\min_{\rank(B)\le k}\|A-B\|=\sigma_{k+1},\\
\min_{\rank(B)\le k}\|A-B\|_F=\sqrt{\sum_{i>k}\sigma_i^2}
\end{gathered}
\]
\(\sigma_r\) is the operator-norm distance to rank \(<r\); for invertible
square \(A\), \(\sigma_{\min}\) is distance to singularity.

\sub{Truncated-SVD inverse:}
\[
x_{\rm TSVD}^{(k)}
=\sum_{i=1}^k\frac{u_i^Ty}{\sigma_i}v_i
\]
Small \(k\) causes bias/oversmoothing; large \(k\) admits
noise-amplifying directions. Choose \(k\) by the error/validation
tradeoff, not just numerical rank.

\sect{PCA (Lecture 18)}
\sub{Center first:} with data as columns \(X=[x_1\cdots x_N]\),
\[
\begin{gathered}
\bar x=\frac1N\sum_i x_i,\qquad
X_c=[x_1-\bar x\ \cdots\ x_N-\bar x],\\
S=\frac1N X_cX_c^T
\end{gathered}
\]
\(S=S^T\succeq0\). If data are stored as rows, use
\(S=X_c^TX_c/N\) and transpose the singular-vector roles.

\sub{Variance direction:} for unit \(v\),
\[
\Var(v^Tx_i)=v^TSv
\]
This is maximized by the top eigenvector \(u_1\) of \(S\), with value
\(\lambda_1\). Subsequent PCs maximize variance subject to orthogonality.
Total variance is
\[
\tr(S)=\sum_i\lambda_i=\frac1N\|X_c\|_F^2
\]

\sub{SVD connection:} if \(X_c=U\Sigma V^T\), then
\[
S=U\frac{\Sigma\Sigma^T}{N}U^T,\qquad
\lambda_i=\frac{\sigma_i^2}{N}
\]
For column-stored data the principal directions are left singular vectors
\(u_i\).

\sub{Projection / reconstruction:} with
\(U_k=[u_1\cdots u_k]\),
\[
z_i=U_k^T(x_i-\bar x),\qquad
\hat x_i=\bar x+U_kz_i
\]
The optimal \(k\)-dimensional affine subspace passes through \(\bar x\)
and uses the top PCs. Its mean-square reconstruction error is
\[
E_k=\frac1N\sum_i\|x_i-\hat x_i\|^2
=\sum_{i>k}\lambda_i
=\frac1N\sum_{i>k}\sigma_i^2
\]
Fraction of variance captured:
\[
\rho_k=\frac{\sum_{i=1}^k\lambda_i}{\sum_i\lambda_i}
=\frac{\sum_{i=1}^k\sigma_i^2}{\sum_i\sigma_i^2}
\]

\sub{Covariance ellipsoid / whitening:} for \(S\succ0\),
\[
(x-\bar x)^TS^{-1}(x-\bar x)\le\alpha
\]
has axes \(u_i\) and lengths \(\sqrt{\alpha\lambda_i}\). The whitened
coordinate \(z=S^{-1/2}(x-\bar x)\) has identity covariance.

\sub{Small singular values:} for
\(Y=[y_1\cdots y_N]\in\R^{m\times N}\), \(N\ge m\),
\[
\min_{\|q\|=1}\frac{\|q^TY\|}{\|Y\|_F}
=\frac{\sigma_{\min}(Y)}{\|Y\|_F}
\]
The minimizing \(q\) is the bottom left singular vector; a small value
reveals \(q^Ty_i\approx0\). Center for affine relations. Spectral gaps
above the noise floor reveal effective dimension/source count.

\sect{Homework 6--7 Patterns}
\sub{Ellipsoid fitting check:} positivity of the fitted eigenvalues proves
the quadratic surface is a bounded ellipsoid. Mean residual near zero can
hide cancellation; report RMS
\(\|r-\one\|/\sqrt N\) and \(\|r-\one\|_\infty\).

\sub{Singular-value proof tools:} for the square matrices in HW7, use
\[
\sigma_{\max}(A)=\max_{\|x\|=1}\|Ax\|,\qquad
\sigma_{\min}(A)=\min_{\|x\|=1}\|Ax\|
\]
then apply ordinary/reverse triangle inequalities. A row norm is
\(\|A^Te_i\|\le\sigma_{\max}(A)\); no analogous lower bound holds for
every row.

\sub{Deconvolution:} smoothing operators have smooth, low-frequency
right singular vectors at large \(\sigma_i\), while high-frequency
directions often have tiny \(\sigma_i\). Direct inversion amplifies noise.
TSVD discards those directions; Tikhonov damps them continuously.

\sub{Tikhonov estimator as a matrix:} for
\[
B_\mu=(A^TA+\mu I)^{-1}A^T
\]
its singular-direction gains are
\[
g_i=\frac{\sigma_i}{\sigma_i^2+\mu},\qquad
\|B_\mu\|=\max_i g_i
\]
The \(g_i\) need not remain sorted in the order of \(\sigma_i\). Also
\[
AB_\mu-I=-U\diag\left(\frac{\mu}{\sigma_i^2+\mu}\right)U^T
\]
for square full SVD, so worst reconstruction bias is governed by the
smallest singular value.

\sect{Probability Basics (Lecture 19)}
\sub{Events / axioms:} sample space \(\Omega\), outcome
\(\omega\in\Omega\), event \(A\subseteq\Omega\). Probability satisfies
\[
\Prob(A)\ge0,\qquad \Prob(\Omega)=1,\qquad
\Prob(A\cup B)=\Prob(A)+\Prob(B)
\]
for disjoint \(A,B\). Hence
\[
\Prob(A^c)=1-\Prob(A),\quad
\Prob(A\cup B)=\Prob(A)+\Prob(B)-\Prob(A\cap B)
\]

\sub{Conditioning / Bayes:}
\[
\Prob(A\mid B)=\frac{\Prob(A\cap B)}{\Prob(B)},\qquad
\Prob(A\mid B)=\frac{\Prob(B\mid A)\Prob(A)}{\Prob(B)}
\]
\(A,B\) are independent iff
\(\Prob(A\cap B)=\Prob(A)\Prob(B)\), equivalently
\(\Prob(A\mid B)=\Prob(A)\).

\sub{Distributions:} discrete PMF \(p(x)=\Prob(X=x)\),
\(\sum_xp(x)=1\). Continuous PDF \(f(x)\ge0\),
\[
\Prob(a\le X\le b)=\int_a^b f(x)\,dx,\qquad \int f=1
\]
CDF \(F(t)=\Prob(X\le t)\) works in both cases; for a continuous variable,
\(F'(t)=f(t)\) where differentiable.

\sub{Common distributions:}
\[
\begin{array}{c|c|c}
X&\E X&\Var X\\ \hline
\operatorname{Bern}(p)&p&p(1-p)\\
\operatorname{Bin}(n,p)&np&np(1-p)\\
\operatorname{Unif}[a,b]&(a+b)/2&(b-a)^2/12\\
\operatorname{Exp}(\lambda)&1/\lambda&1/\lambda^2\\
\mathcal N(\mu,\sigma^2)&\mu&\sigma^2
\end{array}
\]
Also Poisson\((\lambda)\) has mean and variance \(\lambda\).

\sub{Expectation / moments:}
\[
\E[g(X)]=\sum_xg(x)p(x)\quad\text{or}\quad
\int g(x)f(x)\,dx
\]
\[
\E[aX+bY+c]=a\E X+b\E Y+c
\]
Linearity never requires independence. The \(k\)th moment is \(\E X^k\);
the \(k\)th central moment is \(\E[(X-\E X)^k]\).

\sub{Variance:}
\[
\begin{gathered}
\Var X=\E[(X-\mu)^2]=\E[X^2]-\mu^2,\\
\Var(aX+b)=a^2\Var X
\end{gathered}
\]
\[
\Var(X+Y)=\Var X+\Var Y+2\Cov(X,Y)
\]
For independent variables the covariance term is zero.

\sub{Covariance / correlation:}
\[
\Cov(X,Y)=\E[(X-\mu_X)(Y-\mu_Y)]
=\E[XY]-\mu_X\mu_Y
\]
\[
\rho_{XY}=\frac{\Cov(X,Y)}{\sigma_X\sigma_Y}\in[-1,1]
\]
Independence implies zero covariance; zero covariance generally does not
imply independence.

\sub{Random vectors:} for \(x\in\R^n\),
\[
\mu=\E x,\qquad
\Sigma=\Cov(x)=\E[(x-\mu)(x-\mu)^T]\succeq0
\]
Diagonal entries are variances; off-diagonal entries are covariances.
For deterministic \(A,b\),
\[
\E[Ax+b]=A\mu+b,\qquad
\Cov(Ax+b)=A\Sigma A^T
\]
If independent \(x,w\), then
\(\Cov(Ax+w)=A\Sigma_xA^T+\Sigma_w\).

\sub{Sample mean / limit theorems:} for independent identically
distributed \(X_i\) with mean \(\mu\), variance \(\sigma^2\),
\[
\bar X_N=\frac1N\sum_iX_i,\qquad
\E\bar X_N=\mu,\qquad
\Var(\bar X_N)=\frac{\sigma^2}{N}
\]
LLN: \(\bar X_N\to\mu\). CLT: for large \(N\),
\(\bar X_N\approx\mathcal N(\mu,\sigma^2/N)\).

\sub{Gaussian facts:} \(X\sim\mathcal N(\mu,\sigma^2)\) has density
\[
f(x)=\frac1{\sqrt{2\pi\sigma^2}}
\exp\left(-\frac{(x-\mu)^2}{2\sigma^2}\right)
\]
About \(68\%,95\%,99.7\%\) lies within \(1,2,3\) standard deviations.
Affine transformations of jointly Gaussian vectors are jointly Gaussian.
For jointly Gaussian variables, zero covariance implies independence.

\sect{Gaussian Conditioning / Practice Final}
\sub{Joint Gaussian conditioning:} if
\[
\begin{bmatrix}x\\y\end{bmatrix}\sim\mathcal N\left(
\begin{bmatrix}\mu_x\\\mu_y\end{bmatrix},
\begin{bmatrix}\Sigma_{xx}&\Sigma_{xy}\\
\Sigma_{yx}&\Sigma_{yy}\end{bmatrix}\right)
\]
then
\[
\E[x\mid y]=\mu_x+\Sigma_{xy}\Sigma_{yy}^{-1}(y-\mu_y)
\]
\[
\Cov(x\mid y)=
\Sigma_{xx}-\Sigma_{xy}\Sigma_{yy}^{-1}\Sigma_{yx}
\]
The conditional covariance does not depend on the observed value \(y\).

\sub{Linear Gaussian model:} if
\[
x\sim\mathcal N(0,\tau^2I),\quad y=Ax+w,\quad
w\sim\mathcal N(0,\sigma^2I),\quad x\perp w
\]
then
\[
\Cov\begin{bmatrix}x\\y\end{bmatrix}
=
\begin{bmatrix}
\tau^2I&\tau^2A^T\\
\tau^2A&\tau^2AA^T+\sigma^2I
\end{bmatrix}
\]
and the MMSE estimator is
\[
\hat x_{\rm MMSE}
=\tau^2A^T(\tau^2AA^T+\sigma^2I)^{-1}y
\]

\sub{MMSE equals regularized LS:}
\[
\hat x_{\rm reg}=(A^TA+\lambda I)^{-1}A^Ty,\qquad
\boxed{\lambda=\sigma^2/\tau^2}
\]
Use \(B(I+CB)^{-1}=(I+BC)^{-1}B\) or the push-through identity.

\sub{MMSE shrinkage in SVD:} for \(A=U\Sigma V^T\),
\[
\hat x_{\rm MMSE}
=\sum_i\frac{\tau^2\sigma_i}{\tau^2\sigma_i^2+\sigma^2}
(u_i^Ty)v_i
\]
Relative to \(x_{\rm ls}\), component \(i\) is multiplied by
\[
\gamma_i=\frac{\tau^2\sigma_i^2}
{\tau^2\sigma_i^2+\sigma^2}\in(0,1)
\]
Large gain/low noise/weak prior shrinkage gives \(\gamma_i\approx1\);
small gain/high noise/strong small-\(x\) prior gives \(\gamma_i\approx0\).

\sub{Scalar posterior MSE:} for \(y=sx+w\) with \(s>0\),
\[
\E[(x-\hat x_{\rm MMSE})^2\mid y]
=\frac{\tau^2\sigma^2}{\sigma^2+s^2\tau^2}
\]
It tends to \(0\) as \(\sigma\to0\) and to \(\tau^2\) as
\(\sigma\to\infty\).

\sect{Practice Final Problem Patterns}
\sub{Warehouse stacking:} for two warehouses,
\[
A_{\rm ship}=
\begin{bmatrix}I&0\\I&0\\0&I\\0&I\end{bmatrix},\quad
y=\begin{bmatrix}c_1\\c_2\\c_3\\c_4\end{bmatrix},\quad
A_{\rm pen}=[I\ -I]
\]
Use \(\hat A=[A_{\rm ship};\sqrt\mu A_{\rm pen}]\),
\(\hat y=[y;0]\). Uniqueness holds because \(A_{\rm ship}\) has full
column rank. At \(\mu=0\), each warehouse is its pair's midpoint; as
\(\mu\to\infty\), both approach the centroid of all cities. In the
practice-final diagonalization, varying \(\mu\) traces the line segment
between these endpoint solutions.

\sub{Separate regressions:} fitting nonorthogonal feature columns one at a
time generally double-counts shared directions and gives a worse residual
than joint LS. With orthonormal feature columns \(U\), coefficients decouple:
\[
\min_\beta\left\|y-\sum_j\beta_ju_j\right\|^2
\quad\Rightarrow\quad \beta_j=u_j^Ty
\]
The fitted vector matches joint LS, although coordinates differ from those
in the original feature basis.

\sub{Rank-deficient regression:} zero singular values make \(X^TX\)
singular. Use \(X^\dagger y\); it gives the minimum-norm coefficient among
all LS minimizers, while \(XX^\dagger y\) is the unique fitted response.

\sub{Graph Laplacian:} for an undirected weighted graph,
\[
L=D-W,\qquad
x^TLx=\frac12\sum_{i,j}w_{ij}(x_i-x_j)^2\ge0
\]
Thus \(L\succeq0\), \(L\one=0\), and the multiplicity of eigenvalue zero
equals the number of connected components. The graph is connected iff
\(\nullsp(L)=\spn\{\one\}\), equivalently \(\lambda_2(L)>0\).

\end{multicols*}
\end{document}
